 \subsection  {Κατηγορίες ρύπων}


\subsubsection {Αέρινα σωματίδια}



Τα αέρινα σωματίδια (PM) αποτελούν σύνθετο μείγμα οργανικών και ανόργανων χημικών ενώσεων, περιλαμβάνοντας οργανικό άνθρακα (OC) και στοιχειακό άνθρακα (EC), των οποίων οι συγκεντρώσεις διαφέρουν ανάλογα με την εποχή και την γεωγραφική θέση [16,17]. 
Τα εσωτερικά σωματίδια περιλαμβάνουν αέρια σωματίδια διαφορετικών κλασμάτων μεγέθους με βάση  την διάμετρος τους. 
Αν η διάμετρος τους είναι από 2.5 έως 10  μικρόμετρα, ονομάζονται χονδρόκοκκα σωματίδια (coarse  particles) ή  PM10.
Αν η διάμετρος τους είναι κάτω από 2.5 μικρόμετρα, ονομάζονται λεπτά σωματίδια (fine particles) ή  PM2.5.
Και για διάμετρο κάτω των από 1.0 μικρόμετρα, ονομάζονται υπέρλεπτα σωματίδια – (Ultrafine particles ,UFPs) ή PM1, 

Ο διαχωρισμός υπάρχει γιατί τα  PM2.5,λόγω του πολύ μικρού μεγέθους τους, μπορούν να  εισχωρούν στο εσωτερικό των πνευμόνων, στην επιφάνεια των επιθηλιακών κυττάρων των βρογχιολίων και των κυψελίδων, ενώ τα PM1.0 μπορούν να διεισδύσουν σε εσωτερικά όργανα, όπως η καρδιά και ο εγκέφαλος \cite {SANGKHAM2024100603}. 

Τα PM2.5 και PM1.0 μπορούν να προκαλέσουν πνευμονικές λοιμώξεις και να δημιουργήσουν αγγειακές και ενδοθηλιακές δυσλειτουργίες, διαταραχές της μεταβλητότητας της καρδιακής συχνότητας, διαταραχές της πήξης και της καρδιακής αυτόνομης λειτουργίας .Η μακροχρόνια έκθεση σε αέρινα σωματίδια μπορεί να προκαλέσει πλήθος αναπνευστικών και καρδιαγγειακών προβλημάτων, από ερεθισμό ματιών, μύτης, λαιμού και βρόγχων, καθώς και άσθμα, έως ίνωση, ανθράκωση και καρκίνο των πνευμόνων.\cite{su12219045},\cite{GONZALEZMARTIN2021128376}

Τα PM εισέρχονται από το εξωτερικό περιβάλλον, αλλά και παράγονται από τον εσωτερικό χώρο. 
Οι κύριες εσωτερικές πηγές PM, για τα λεπτά σωματίδια PM2.5 περιλαμβάνουν το κάπνισμα, το μαγείρεμα, ιδιαίτερα με χρήση κηροζίνης και καυσίμων βιομάζας, ξυλόσομπες και φούρνους, τη χρήση θυμιάματος και κεριών.
Για τα χονδρόκοκκα σωματίδια PM10,περιλαμβάνονται δραστηριότητες καθαρισμού, παρουσία κατοικίδιων και ανθρώπινες κινήσεις για τα χονδρόκοκκο σωματίδια [20–26]\cite{ijerph17238972}.


\subsubsection{Πτητικές χημικές ουσίες (Volatile Organic Compounds, VOCs)}

Οι πτητικές οργανικές ενώσεις (Volatile Organic Compounds, VOCs) αναγνωρίζονται ως αέρια που περιέχουν ποικίλα χημικά συστατικά και εκπέμπονται από υγρά ή στερεά [34]. Οι συγκεντρώσεις των VOCs σε εσωτερικούς χώρους είναι τουλάχιστον 10 φορές υψηλότερες από τους εξωτερικούς χώρους, ανεξαρτήτως της τοποθεσίας του κτιρίου [34,35].
Δεδομένου ότι οι VOCs είναι οργανικές ενώσεις με χαμηλό σημείο βρασμού (Tb) και εξατμίζονται εύκολα σε θερμοκρασία δωματίου, ο Παγκόσμιος Οργανισμός Υγείας (WHO) τις ταξινόμησε σε τέσσερις ομάδες: 
\begin {enumerate}
    \item Πολύ πτητικές οργανικές ενώσεις (VVOCs) με θερμοκρασία βρασμού 50–100 \degree C.
    \item Πτητικές οργανικές ενώσεις (VOCs) με 100 °C < Tb < 240 \degree C.
    \item Ήμι-πτητικές ανόργανες ενώσεις (SVOCs) με 240 \degree C < Tb < 380 \degree C.
    \item Σωματιδιακή οργανική ύλη (POM) με Tb > 380 \degree °C.
\end{enumerate}

Το σύνολο των VOCs σε ένα χώρο ονομάζονται Total Volatile Organic Compounds(TVOCs)

Η έκθεση στα VOCs που απελευθερώνονται από καταναλωτικά προϊόντα γίνεται κυρίως μέσω τριών οδών: εισπνοή, κατάποση ή δερματική επαφή. Οι περισσότεροι άνθρωποι δεν επηρεάζονται σοβαρά από βραχυχρόνια έκθεση σε χαμηλές συγκεντρώσεις VOCs, αλλά σε περιπτώσεις μακροχρόνιας έκθεσης, ορισμένες VOCs θεωρούνται επικίνδυνες για την ανθρώπινη υγεία, με πιθανή καρκινογόνο δράση [43]. Όσον αφορά τις SVOCs, η δερματική απορρόφηση απευθείας από τον αέρα συμβάλλει περισσότερο στην έκθεση σε σύγκριση με την εισπνοή [44,45].\cite{ijerph17082927}

Πολλές μελέτες έχουν δείξει ότι οι συγκεντρώσεις πολλών VOCs σε εσωτερικούς χώρους είναι σημαντικά υψηλότερες από εκείνες των εξωτερικών χώρων . Οι κύριες πηγές εσωτερικών VOCs περιλαμβάνουν υλικά κατασκευής με υψηλές εκπομπές, έπιπλα, αερολύματα, φυτοφάρμακα, υφάσματα, καθώς και εξοπλισμό γραφείου όπως φωτοτυπικά μηχανήματα και εκτυπωτές λέιζερ . Η Υπηρεσία Προστασίας Περιβάλλοντος των ΗΠΑ (US EPA) έχει εκδώσει κατάλογο επικίνδυνων ατμοσφαιρικών ρύπων, που περιλαμβάνει συνολικά 187 VOCs \cite{su12219045},\cite{USAEPA2025}. Επιπλέον, το πρότυπο ANSI/ASHRAE 62.1-2016 παρέχει τις Επίπεδα Αναφοράς Έκθεσης (Reference Exposure Levels, RELs) για 32 συγκεκριμένες VOCs εσωτερικού χώρου για τον γενικό πληθυσμό [74].

Οι πιο συνηθισμένες εσωτερικές VOCs—όπως το βενζόλιο, η αιθυλένιο, η φορμαλδεΰδη, η μεθυλενιοχλωρίδιο, η τετραχλωροαιθυλένιο, η τολουόλη, η ξυλόλη και η 1,3-βουταδιένιο—έχουν αποδειχθεί ότι αποτελούν καρκινογόνες, ερεθιστικές και τοξικές ουσίες για τον άνθρωπο . , αναπνευστικά προβλήματα, ενώ ορισμένα μπορεί να προκαλέσουν καρκίνο [67].  

Τα TVOCs χρησιμοποιούνται ως μέτρο της συνολικής συγκέντρωσης VOCs σε εσωτερικούς χώρους. Η οξεία έκθεση σε TVOCs μπορεί να προκαλέσει ερεθισμό ματιών, μύτης και λαιμού, πονοκεφάλους, απώλεια συντονισμού και ναυτία, βλάβες στο ήπαρ, τα νεφρά και το κεντρικό νευρικό σύστημα. Τα άτομα με άσθμα, τα μικρά παιδιά και οι ηλικιωμένοι είναι ιδιαίτερα ευάλωτοι στις επιπτώσεις της έκθεσης σε TVOCs

Πολλές VOCs ταξινομούνται ως γνωστά ή πιθανά καρκινογόνα (ιδιαίτερα το βενζόλιο, γνωστό καρκινογόνο για τον άνθρωπο κυρίως συσχετισμένο με λευχαιμία), ερεθιστικά και τοξικά, ενώ η μέτρηση των TVOCs μπορεί να υποεκτιμά τους κινδύνους που σχετίζονται με μεμονωμένες ενώσεις [5,70]. Οι περισσότερες VOCs έχουν επίσης αναφερθεί ως σημαντικοί παράγοντες κινδύνου για άσθμα [73,74,90–92], με τη μεγαλύτερη συσχέτιση για το βενζόλιο, ακολουθούμενο από αιθυλοβενζόλιο και τολουόλη \cite{ijerph17238972},\cite{su12219045}.


\subsubsection{Λίστα χαρακτηριστικών πτητικών χημικές ουσίες}

Οι παρακάτω ουσίες περιγράφονται από τους \cite{GONZALEZMARTIN2021128376}:

\paragraph{Βενζόλιο (Benzene)}

Το βενζόλιο είναι μια πτητική οργανική ένωση ευρείας χρήσης στη βιομηχανία. Οι εσωτερικές πηγές βενζολίου περιλαμβάνουν τον περιβαλλοντικό καπνό καπνού (ETS), συσκευές καύσης και δομικά υλικά, όπως πολυμερή έπιπλα, χαλιά, χρώματα, διαλύτες, διάφορα πλαστικά και ξύλινα έπιπλα. Οι εκπομπές βενζολίου (και άλλων ρύπων) από τα υλικά αυτά μειώνονται με την πάροδο του χρόνου· ως εκ τούτου, τα επίπεδα είναι σημαντικά υψηλότερα σε νέα ή πρόσφατα ανακαινισμένα κτίρια.

Οι κίνδυνοι για την υγεία από οξεία δηλητηρίαση θεωρούνται χαμηλοί (π.χ. πονοκέφαλοι μετά από έκθεση 5 ωρών σε συγκεντρώσεις 150–500 mg/m\textsuperscript{3}), ωστόσο το βενζόλιο είναι αποδεδειγμένα καρκινογόνο και γενοτοξικό σε συνθήκες χρόνιας έκθεσης. Επιπλέον, μπορεί να προκαλέσει αιματολογικές παθήσεις καθώς και νευρολογικές και αναπαραγωγικές διαταραχές. Δεν συνιστάται ασφαλές επίπεδο έκθεσης στο βενζόλιο.

Οι υψηλότερες συγκεντρώσεις βενζολίου έχουν καταγραφεί σε εμπορικούς χώρους (2,5–48 μg/m\textsuperscript{3}), ακολουθούμενες από γραφεία (1,4–5,5 μg/m\textsuperscript{3}), κατοικίες (0,7–4,4 μg/m\textsuperscript{3}) και σχολεία (0,5–3 μg/m\textsuperscript{3}),

\paragraph{τολουόλιο (Toluene)}

Το τολουόλιο, η οποία απαντάται στη βενζίνη, τα χρώματα, τις ρητίνες, τα συγκολλητικά, τα καλλυντικά και τις επιστρώσεις, αποτελεί έναν από τους συχνότερους ρύπους του εσωτερικού αέρα και τη σημαντικότερη συνιστώσα της ομάδας BTEX (βενζόλιο, τολουόλιο, αιθυλοβενζόλιο και ξυλόλιο). Οι κύριες πηγές Τολουολίου σε εσωτερικούς χώρους περιλαμβάνουν τη διείσδυση από τον εξωτερικό αέρα, τον καπνό καπνού, τις συσκευές καύσης και ένα ευρύ φάσμα οικιακών προϊόντων 

Το τολουόλιο δεν έχει επιβεβαιωθεί ως καρκινογόνος ουσία και, ως εκ τούτου, οι μακροχρόνιες καρκινογενετικές επιπτώσεις της δεν έχουν τεκμηριωθεί επαρκώς. Παρ’ όλα αυτά, η έκθεση Το τολουόλιο έχει συσχετιστεί με αυξημένο κίνδυνο εμφάνισης άσθματος και άλλων αναπνευστικών παθήσεων, ενώ τόσο η βραχυχρόνια όσο και η μακροχρόνια έκθεση μπορεί να επηρεάσει το κεντρικό νευρικό σύστημα.

Οι υψηλότερες συγκεντρώσεις τολουόλης έχουν καταγραφεί σε εμπορικά κέντρα (15–164 μg/m\textsuperscript{3}), ακολουθούμενες από χώρους γραφείων (6–32 μg/m\textsuperscript{3}) και κατοικίες (3–20 μg/m\textsuperscript{3}). Η παρουσία τολουόλης σε σχολεία είναι περιορισμένη, με τυπικές συγκεντρώσεις της τάξης των 1,8 μg/m\textsuperscript{3}. Οι συγκεντρώσεις είναι συνήθως υψηλότερες σε νέα ή πρόσφατα ανακαινισμένα κτίρια, γεγονός που αποδίδεται στην εκπομπή ρύπων από τα καινούργια εγκαταστημένα υλικά.

\paragraph{Αιθυλοβενζόλιο (Ethylbenzene)}


Το αιθυλοβενζόλιο είναι ένα αρωματικό, άχρωμο υγρό σε θερμοκρασία περιβάλλοντος. Οι πηγές ρύπανσης περιλαμβάνουν πλαστικά, χρώματα, συγκολλητικά και άλλα προϊόντα στα οποία το αιθυλοβενζόλιο χρησιμοποιείται ως διαλύτης κατά τη διαδικασία παραγωγής. Επιπλέον, αιθυλοβενζόλιο εκπέμπεται και κατά τις διεργασίες καύσης.

Οι οξείες επιδράσεις σε σχετικά υψηλές συγκεντρώσεις περιλαμβάνουν ερεθισμό των οφθαλμών και του φάρυγγα, καθώς και ζάλη. Μέχρι σήμερα δεν έχουν τεκμηριωθεί μακροχρόνιες επιπτώσεις στην ανθρώπινη υγεία. Ωστόσο, Το αιθυλοβενζόλιο ε έχει ταξινομηθεί ως δυνητικά καρκινογόνος για τον άνθρωπο από τον Διεθνή Οργανισμό Έρευνας για τον Καρκίνο.

Συγκεντρώσεις Αιθυλοβενζολίου έως και 16 μg/m\textsuperscript{3} έχουν ανιχνευθεί σε χώρους γραφείων και εμπορικά κέντρα. Αντιθέτως, οι κατοικίες και τα σχολεία εμφανίζουν μέγιστες συγκεντρώσεις της τάξης των 8 και 0,2 μg/m\textsuperscript{3}, αντίστοιχα.

\paragraph{Ξυλόλιο (Xylene)}

Το ξυλόλιο απαντώνται σε τρία διαφορετικά ισομερή: ορθό-, μετά- και πάρα- (o-, m- και p-, αντίστοιχα). Οι πηγές ξυλολίου σε εσωτερικούς χώρους περιλαμβάνουν χρώματα, συγκολλητικές ουσίες, βαφές, πολυμερή, καθαριστικά προϊόντα και φαρμακευτικά σκευάσματα.

Η οξεία έκθεση στο ξυλόλιο μπορεί να προκαλέσει ερεθισμό των οφθαλμών και του φάρυγγα, κεφαλαλγία και ναυτία. Η μακροχρόνια έκθεση έχει συσχετιστεί με διαταραχές του αναπνευστικού, του γαστρεντερικού και του κεντρικού νευρικού συστήματος, καθώς και με επιπτώσεις στους πνεύμονες, τους νεφρούς, την καρδιά και το αναπαραγωγικό σύστημα. Τα ισομερή ξυλολίου θεωρούνται ύποπτες καρκινογόνες ουσίες και έχουν συσχετιστεί με αυξημένο κίνδυνο εμφάνισης λευχαιμίας, μη-Hodgkin λεμφώματος και καρκίνου του παχέος εντέρου/ορθού.

Οι εμπορικοί χώροι συγκαταλέγονται μεταξύ των πλέον επιβαρυμένων περιβαλλόντων, με συγκεντρώσεις που κυμαίνονται από 1,3 έως 74 μg/m\textsuperscript{3}, ενώ σε χώρους γραφείων έχουν καταγραφεί επίπεδα 2,2–16 μg/m\textsuperscript{3}. Αντιθέτως, οι μέγιστες συγκεντρώσεις ξυλολιών που έχουν μετρηθεί σε κατοικίες και σχολεία είναι σημαντικά χαμηλότερες, της τάξης των 3,1 και 0,3 μg/m\textsuperscript{3} αντίστοιχα.

\paragraph{Ναφθαλίνη (Naphthalene)}

Η ναφθαλίνη είναι ένα λευκό κρυσταλλικό στερεό, το οποίο χρησιμοποιείται ευρέως ως πρώτη ύλη στη χημική βιομηχανία. Πηγές ναφθαλίνης σε εσωτερικούς χώρους περιλαμβάνουν τον καπνό του τσιγάρου, ελαττωματικές συσκευές καύσης, ζιζανιοκτόνα, ελαστικά υλικά και εντομοαπωθητικά ,συνήθως μέσω επιμόλυνσης των ρούχων, καθώς και τον εξωτερικό αέρα.

Η εισπνοή αποτελεί τον κύριο μηχανισμό πρόσληψης της ναφθαλίνης στον άνθρωπο, αν και έχουν αναφερθεί περιπτώσεις τυχαίας κατάποσης σφαιριδίων ναφθαλίνης. Η οξεία δηλητηρίαση από ναφθαλίνη μπορεί να προκαλέσει αιμολυτική αναιμία και καταρράκτη. Δεν έχουν αναφερθεί μακροχρόνιες δυσμενείς επιπτώσεις στην υγεία ούτε τεκμηριωμένα στοιχεία καρκινογένεσης ή γονοτοξικότητας στον άνθρωπο.

Συγκεντρώσεις ναφθαλίνης που κυμαίνονται από 3 έως 26 μg/m\textsuperscript{3} έχουν καταγραφεί σε κατοικίες.


\paragraph{Φορμαλδεΰδη (Formaldehyde)}

Η φορμαλδεΰδη (HCHO) είναι ένα ιδιαίτερα δραστικό αέριο που εκλύεται κατά την ατελή αντίδραση υδρογονανθράκων. Η φορμαλδεΰδη παράγεται επίσης από την οξείδωση άλλων VOCs με όζον ή ακτινοβολία. Αυτή η χημική ουσία χρησιμοποιείται ευρέως σε ρητίνες, κόλλες, χρώματα, προϊόντα χαρτιού, καλλυντικά, ηλεκτρονικό εξοπλισμό, καθαριστικά και υφάσματα. Η φορμαλδεΰδη είναι παρούσα σε δομικά υλικά, όπως οι μονωτικές αφρώδεις ύλες ξύλινων προϊόντων που χρησιμοποιούνται σε δάπεδα ή έπιπλα. Ωστόσο, οι εκπομπές από αυτά τα υλικά (π.χ. κόντρα πλακέ, μοριοσανίδες ή ίνες ξύλου) συνήθως μειώνονται μέσα σε μερικές εβδομάδες.

Η φορμαλδεΰδη απορροφάται γρήγορα από το αναπνευστικό ή το γαστρεντερικό σύστημα. Οι οξείες επιδράσεις της έκθεσης περιλαμβάνουν οσμή, ερεθισμό, πονοκέφαλο και έκζεμα. Η φορμαλδεΰδη είναι γνωστό καρκινογόνο και γονοτοξικό χημικό, της οποίας η μακροχρόνια έκθεση μπορεί να προκαλέσει καρκίνο του ρινοφάρυγγα και μυελοειδή λευχαιμία.

Τα επίπεδα έκθεσης σε σπίτια, γραφεία και σχολεία είναι συγκρίσιμα (7,7–30, 8–17, 9–17 μg/m\textsuperscript{3}).


\paragraph{Τριχλωροαιθυλένιο (Trichloroethylene)}

Το τριχλωροαιθυλένιο (Trichloroethylene, TCE) είναι μια πτητική χημική ουσία που χρησιμοποιείται ευρέως ως διαλύτης και καθαριστικό μέσο (απολίπανση μετάλλων με ατμούς). Πηγές TCE σε εσωτερικούς χώρους περιλαμβάνουν προϊόντα στα οποία χρησιμοποιείται ως διαλύτης, όπως λιπαντικά, βερνίκια, αφαιρετικά χρωμάτων, κόλλες και υγρά διόρθωσης για γραφομηχανές. Το TCE μπορεί επίσης να υπάρχει σε ορισμένα οικιακά προϊόντα λεύκανσης και άλλα καθαριστικά. Επιπλέον, το TCE μπορεί να εντοπίζεται σε μολυσμένα ύδατα και εδάφη, τα οποία ενδέχεται να συμβάλλουν έμμεσα στα επίπεδα συγκέντρωσής του σε εσωτερικούς χώρου.

Η ανθρώπινη έκθεση σε αυτήν τη VOC πραγματοποιείται μέσω δερματικής απορρόφησης, κατάποσης και εισπνοής. Η έκθεση στο TCE είναι συχνά διαλείπουσα, καθώς τα προϊόντα που περιέχουν αυτή την ουσία δεν χρησιμοποιούνται συνεχώς. Η οξεία έκθεση στο TCE επηρεάζει το κεντρικό νευρικό σύστημα σε συγκεντρώσεις περίπου 270 mg/m\textsuperscript{3}, προκαλώντας μείωση της αισθητηριακής ικανότητας. Το TCE είναι καρκινογόνος ουσία σε περίπτωση μακροχρόνιας έκθεσης, ενώ η χρόνια έκθεση μπορεί να προκαλέσει καρκίνο του ήπατος, των νεφρών και των χοληφόρων οδών, καθώς και λέμφωμα non-Hodgkin. Για τον λόγο αυτό, δεν συνιστάται ασφαλές όριο έκθεσης λόγω της καρκινογόνου δράσης του.

Συγκεντρώσεις TCE που κυμαίνονται από 0,3 έως 0,6 μg/m\textsuperscript{3} καταγράφονται συνήθως σε κατοικίες, ενώ σε εμπορικούς χώρους κυμαίνονται από 0,6 έως 4,7 μg/m\textsuperscript{3}.

\paragraph{Α-Πινένιο (Alpha-pinene )}

Το α-πινένιο (Άλφα-πινένιο) είναι ένα υγρό τερπένιο που απαντάται φυσικά στα φυτά και αποτελεί κοινό συστατικό των αιθέριων ελαίων. Το  α-πινένιο χρησιμοποιείται ως διαλύτης σε ορισμένα χρώματα και υλικά στεγανοποίησης και μπορεί να βρεθεί σε προϊόντα αρωματοποιίας, αποσμητικά και καθαριστικά προϊόντα. Εκπομπές  α-πινενίου μπορούν επίσης να προέρχονται από ξύλινα υλικά, όπως έπιπλα ή δάπεδα, ιδιαίτερα όταν αυτά είναι κατασκευασμένα από ξύλο πεύκου.

Τα διαθέσιμα δεδομένα σχετικά με τους κινδύνους για την υγεία από το α-πινένιο είναι περιορισμένα. Αν και η οξεία έκθεση σε υψηλές συγκεντρώσεις μπορεί να προκαλέσει ερεθισμό και φλεγμονή, δεν έχουν αναφερθεί μέχρι σήμερα κίνδυνοι από χρόνια έκθεση. Ωστόσο, ενδέχεται να προκύψουν κίνδυνοι για την υγεία όταν συνυπάρχει όζον ή άλλες δραστικές χημικές ρίζες, λόγω της υψηλής αντιδραστικότητας του  α-πινενίου και άλλων ακόρεστων μορίων. Τα παραπροϊόντα αυτών των αντιδράσεων μπορούν να προκαλέσουν ερεθισμό των οφθαλμών και των ανώτερων αεραγωγών σε μεγαλύτερο βαθμό και για μεγαλύτερη διάρκεια σε σύγκριση με το ίδιο το α-πινένιο.

Οι υψηλότερες συγκεντρώσεις πινενίου καταγράφονται σε κατοικίες (11–32 μg/m\textsuperscript{3}) και βιβλιοθήκες (10–30 μg/m\textsuperscript{3}), ενώ σε γραφεία και σχολεία παρατηρούνται συγκεντρώσεις έως 6,3 και 1,5 μg/m\textsuperscript{3} αντίστοιχα.

\paragraph {Λιμονένιο (Limonene)}

Το λιμονένιο είναι ένα φυσικό τερπένιο που υπάρχει σε δύο ισομερείς μορφές (d- και l-). Χρησιμοποιείται σε μια μεγάλη ποικιλία οικιακών προϊόντων, όπως καθαριστικά, ρητίνες, αρωματικά χώρου, αποσμητικά, αρώματα και σαμπουάν. Το τερπένιο αυτό βρίσκεται επίσης στα τρόφιμα ως πρόσθετο λόγω της κιτρικής του οσμής και γεύσης.

Η έκθεση στο λιμονένιο γίνεται κυρίως μέσω εισπνοής ή κατάποσης. Ο συγκεκριμένος ρύπος έχει χαμηλή οξεία τοξικότητα και μέχρι στιγμής δεν έχουν περιγραφεί συμπτώματα πέρα από ερεθισμό των ματιών ή του δέρματος. Δεν έχουν μελετηθεί οι μακροπρόθεσμες επιπτώσεις στην υγεία από το λιμονένιο και δεν υπάρχει ένδειξη καρκινογένεσης ή γονοτοξικότητας. Οι κίνδυνοι για την υγεία συνδέονται κυρίως με την αντίδραση του λιμονενίου με το όζον ή άλλες δραστικές ρίζες στο περιβάλλον, που δημιουργούν παραπροϊόντα όπως αλδεΰδες, καρβοξυλικά οξέα και υπεροξείδια, τα οποία ευθύνονται για ερεθισμό και δυσάρεστη οσμή ακόμα και σε χαμηλές συγκεντρώσεις.

'
\subsubsection{Μονοξείδιο του άνθρακα (CO)}
Το μονοξείδιο του άνθρακα στον εσωτερικό αέρα παράγεται κυρίως από διεργασίες καύσης, όπως το μαγείρεμα ή η θέρμανση. Επιπλέον, το CO μπορεί να εισέλθει σε εσωτερικούς χώρους μέσω διείσδυσης από τον εξωτερικό αέρα . Οι σημαντικές πηγές εκπομπής CO σε εσωτερικά περιβάλλοντα περιλαμβάνουν τα θερμαντικά σώματα κηροζίνης ή αερίου χωρίς εξαερισμό, διαρροές από καπνοδόχους και φούρνους, αναστροφή καπνού από φούρνους, θερμοσίφωνες αερίου, ξυλόσομπες και τζάκια,
κουζίνες αερίου, γεννήτριες και άλλους εξοπλισμούς που λειτουργεί με βενζίνη και καπνό τσιγάρου

Η μέση συγκέντρωση CO σε ένα κτίριο χωρίς κουζίνες αερίου κυμαίνεται περίπου από 0,5–5 ppm, ενώ σε περιοχές κοντά σε κουζίνες αερίου φτάνει από 5 έως 15 ppm και σε ορισμένες περιπτώσεις έως 30 ppm ή περισσότερο. Η έκθεση σε CO μπορεί να προκαλέσει σε χαμηλές συγκεντρώσεις, επιδράσεις στο καρδιαγγειακό σύστημα και στις νευρο-συμπεριφορικές λειτουργίες και σε υψηλές συγκεντρώσεις, απώλεια συνείδησης ή θάνατο \cite{Raub20001},\cite{ijerph17082927}.


\subsubsection{Διοξείδιο του άνθρακα (CO2)}
Το CO2 είναι ένα άχρωμο και άοσμο αέριο, γνωστό συστατικό της ατμόσφαιρας της Γης και επίσης σημαντικό ανθρώπινο μεταβολικό προϊόν. Η μέση συγκέντρωσή του στον εξωτερικό αέρα είναι περίπου 400 ppm, που προέρχεται κυρίως από την καύση ορυκτών καυσίμων . Πρόσφατα, το επίπεδο CO₂ σε εσωτερικούς χώρους χρησιμοποιείται ως δείκτης για την αξιολόγηση της ποιότητας του αέρα (IAQ) και για τον έλεγχο του εξαερισμού . Σύμφωνα με το πρότυπο ASHRAE, προτείνεται οι συγκεντρώσεις CO2 σε εσωτερικούς χώρους να παραμένουν κάτω από 700 ppm για τη διασφάλιση της ανθρώπινης υγείας. Έχει διαπιστωθεί ότι η έκθεση σε συγκέντρωση CO2 3000 ppm αυξάνει την ένταση των πονοκεφάλων, την υπνηλία, την κόπωση και τις δυσκολίες συγκέντρωσης \cite{ijerph17082927}.



\subsubsection{Διοξείδιο του θείου (SO2)}

Το διοξείδιο του θείου (SO2) είναι το πιο κοινό αέριο στην ομάδα των θειωδών οξειδίων (SOx) που βρίσκονται στην ατμόσφαιρα. Το SO2 παράγεται κυρίως από διαδικασίες καύσης ορυκτών καυσίμων και συνδέεται με αερολύματα και σωματίδια (PM) για να σχηματίσει ένα πολύπλοκο σύνολο ατμοσφαιρικών ρύπων [59].

Οι εσωτερικές πηγές εκπομπής SO2 περιλαμβάνουν θερμαντικά σώματα αερίου με εξαερισμό, φούρνους πετρελαίου, καπνό τσιγάρου, θερμαντικά σώματα κηροζίνης, σόμπες ή φούρνους με ξύλο ή άνθρακα.
Επιπλέον, ο εξωτερικός αέρας θεωρείται επίσης κύρια πηγή για SO2 εσωτερικού  αέρα. Οι συγκεντρώσεις SO₂ σε εσωτερικούς χώρους είναι συχνά χαμηλότερες από τις αντίστοιχες εξωτερικές. Οι εσωτερικές εκπομπές SO₂ είναι συνήθως μικρές, λόγω της επαναδραστηριοποίησης και της εύκολης απορρόφησής του από επιφάνειες μέσα στο κτίριο. Είναι γνωστό ότι η ωριαία συγκέντρωση SO2 σε κτίρια συχνά είναι κάτω από 20 ppb.Η έκθεση του ανθρώπου στο SO₂ γίνεται αποκλειστικά μέσω εισπνοής, και μπορεί να προκαλέσει δυσλειτουργία του αναπνευστικού συστήματος.

\cite{ijerph17082927}


\subsubsection{Οξείδια του αζώτου (NOx)}

%bres kalo orismo pou sumperilambanei tin exoteriki kausi ton autokiniton

Το μονοξείδιο του αζώτου (NO) και το διοξείδιο του αζώτου (NO₂) είναι οι πιο συνηθισμένες και σημαντικές ενώσεις οξειδίων του αζώτου. Το NO παράγεται σε μεγαλύτερο βαθμό και οξειδώνεται περαιτέρω σε NO₂ παρουσία οξυγόνου, όζοντος ή VOCs.
Οι πιο συχνοί εσωτερικοί  ρύποι NOx προέρχονται από συσκευές αερίου, όπως κουζίνες, φούρνους ή θερμοσίφωνες. Το παθητικό κάπνισμα και τα τζάκια μπορούν επίσης να συμβάλλουν στις εκπομπές NOx. Η εισχώρηση αέρα από το εξωτερικό επηρεάζει σημαντικά τα εσωτερικά επίπεδα, ιδιαίτερα κοντά σε δρόμους ή περιοχές υψηλής βιομηχανικής πυκνότητας.

Το εισπνεόμενο NO2 μετατρέπεται σε νιτρικό οξύ στους πνεύμονες, προκαλώντας βλάβες στα κύτταρα και στο ανοσοποιητικό σύστημα. Η εισπνοή μπορεί να προκαλέσει αναπνευστική δυσλειτουργία όταν υπάρχουν προϋπάρχουσες παθήσεις (όριο προστασίας για άτομα με άσθμα: 200 μg/m\textsuperscript{3} μέσος όρος 1 ώρας, ≈2000 μg/m\textsuperscript{3} για φυσιολογικά άτομα). Το NO2 μπορεί επίσης να προκαλέσει ερεθισμό των ματιών από άμεση επαφή \cite{GONZALEZMARTIN2021128376}.


\subsubsection{Όζον (Ozone)}

Το όζον είναι ένας ισχυρός οξειδωτικός παράγοντας, που παράγεται κυρίως μέσω φωτοχημικών αντιδράσεων του O2, των NOx και των VOCs στην ατμόσφαιρα. Ωστόσο, δεν μπορεί να χρησιμοποιηθεί για την εξουδετέρωση άλλων χημικών ρύπων στους εσωτερικούς χώρους, λόγω της αργής αντίδρασης του με τους περισσότερους ατμοσφαιρικούς ρύπου \cite{https://doi.org/10.1034/j.1600-0668.2000.010004269.x}.

Το όζον μπορεί να αντιδράσει γρήγορα με αρκετούς εσωτερικούς ρύπους, αλλά τα προϊόντα των αντιδράσεων αυτών μπορούν να προκαλέσουν ερεθισμούς στον άνθρωπο και φθορά στα υλικά. Οι κύριες πηγές εσωτερικού όζοντος προέρχονται από την εξωτερική ατμόσφαιρα και τη λειτουργία ηλεκτρικών συσκευών. Οι συσκευές που συχνά εκπέμπουν όζον στους εσωτερικούς χώρους περιλαμβάνουν φωτοαντιγραφικά μηχανήματα, συσκευές απολύμανσης, καθαρισμού αέρα και άλλες γραφειακές συσκευές. Οι μηχανισμοί εκπομπής όζοντος από αυτές τις συσκευές χωρίζονται στην Κορώνα εκφόρτιση (Corona discharge) και  Φωτοχημικοί μηχανισμοί.

Έχει αποδειχτεί ότι τα επίπεδα του όζοντος στους εσωτερικούς χώρους εξαρτώνται από τα επίπεδα όζοντος στο εξωτερικό περιβάλλον,τα ποσοστά εκπομπής στο εσωτερικό,τα ποσοστά ανταλλαγής αέρα,τα ποσοστά απομάκρυνσης από επιφάνειες και τις αντιδράσεις με άλλες χημικές ουσίες στον αέρα. Συνήθως, τα επίπεδα όζοντος στους εσωτερικούς χώρους κυμαίνονται μεταξύ 20% και 80% των επιπέδων του εξωτερικού όζοντος, ανάλογα με το ποσοστό ανταλλαγής αέρα . Ο άνθρωπος εκτίθεται κυρίως μέσω της εισπνοής, αλλά η επαφή με το δέρμα αναγνωρίζεται επίσης ως διαδρομή έκθεσης 
\cite{ijerph17082927}.

\subsubsection{Βιολογικοί Ρύποι (Biological Pollutants)}


Οι βιολογικοί ρύποι στους εσωτερικούς χώρους περιλαμβάνουν βιολογικούς αλλεργιογόνους παράγοντες, από τις τρίχες και το σάλιο ζώων, την σκόνη σπιτιού, κατσαρίδες, ακάρεα και γύρη) και μικροοργανισμούς (ιούς, μύκητες και βακτήρια).

Οι βιολογικοί αλλεργιογόνοι παράγοντες, γνωστοί και ως αντιγόνα, προέρχονται από διάφορα έντομα, ζώα, ακάρεα, φυτά ή μύκητες, και προκαλούν αλλεργική αντίδραση όταν αλληλεπιδρούν με τα ειδικά αντισώματα ανοσοσφαιρίνη E (lgE). Οι κύριες πηγές αλλεργιογόνων στους εσωτερικούς χώρους περιλαμβάνουν κατοικίδια με τρίχωμα (π.χ. σκύλους και γάτες), ακάρεα σπιτιού, μούχλα, φυτά, κατσαρίδες και τρωκτικά , ενώ υπάρχουν και εξωτερικές πηγές.

Οι ιοί και τα βακτήρια προέρχονται συχνά από ανθρώπους και ζώα ή μεταφέρονται από αυτά. Έχει αποδειχθεί ότι η έκθεση σε βιολογικούς αλλεργιογόνους παράγοντες μπορεί να προκαλέσει ευαισθητοποίηση, αναπνευστικές λοιμώξεις, αλλεργικές αναπνευστικές παθήσεις και συριγμό, ενώ η έκθεση σε βακτήρια και ιούς στους εσωτερικούς χώρους πιθανόν να προκαλέσει τόσο μολυσματικές όσο και μη μολυσματικές αρνητικές επιπτώσεις στην υγεία \cite{ijerph17082927}.

\subsubsection{Ραδόνιο (Radon)}

Οι κύριες πηγές ραδονίου στους εσωτερικούς χώρους περιλαμβάνουν τα οικοδομικά υλικά, τα αέρια του εδάφους και το νερό της βρύσης. Καθώς το έδαφος περιέχει ραδόνιο σε ίχνη, το ραδόνιο είναι πιθανό να αποτελεί ένα από τα συστατικά του αερίου που γεμίζει τους πόρους του εδάφους. Όσον αφορά τις εκπομπές ραδονίου από οικοδομικά υλικά, όλα τα υλικά που περιέχουν ίχνη ραδονίου μπορούν να απελευθερώσουν ραδόνιο. Μεταξύ των οικοδομικών υλικών, τα δομικά υλικά όπως πέτρα, σκυρόδεμα και τούβλο αποτελούν τις κύριες πηγές εκπομπής ραδονίου στους εσωτερικούς χώρους, δεδομένου ότι χρησιμοποιούνται τόνοι αυτών των υλικών στην κατασκευή κτιρίων.

Το ραδόνιο στους εσωτερικούς χώρους μπορεί επίσης να απελευθερωθεί μέσω της χρήσης νερού από υπόγειες πηγές που περιέχουν γρανίτη ή άλλα πετρώματα με ραδόνιο, και τέτοιες πηγές νερού συχνά περιέχουν συγκεντρώσεις ραδονίου άνω των 10.000 pCi/L. Τέλος, ο εξωτερικός αέρας θεωρείται επίσης πηγή εισόδου ραδονίου στους εσωτερικούς χώρους . Η ανθρώπινη έκθεση στο ραδόνιο μέσα σε κτίρια πραγματοποιείται κυρίως μέσω της διάχυσης αερίων από το έδαφος που βρίσκεται κάτω από τα κτίρια . Επιδημιολογικές μελέτες έχουν δείξει ότι το ραδόνιο στους εσωτερικούς χώρους μπορεί να αυξήσει τον κίνδυνο καρκίνου του πνεύμονα κατά 3% έως 14%, ανάλογα με το μέσο επίπεδο ραδονίου \cite{ijerph17082927}.

\subsubsection{Πολυκυκλικοί Αρωματικοί Υδρογονάνθρακες Polycyclic Aromatic Hydrocarbon (PAH)}


Οι PAHs είναι προϊόντα ατελούς καύσης που παράγονται μέσω της καύσης ξύλου, άνθρακα, πετρελαίου και αερίου, του καπνίσματος, της καύσης αποβλήτων, της βιομηχανικής παραγωγής ενέργειας και των εκπομπών οχημάτων. Οι PAHs χαμηλού μοριακού βάρους, με 2 ή 3 δακτυλίους, εκπέμπονται στη αέρια φάση, ενώ οι PAHs υψηλού μοριακού βάρους, με 5 ή περισσότερους δακτυλίους, εκπέμπονται κυρίως στη σωματιδιακή φάση.

 Οι PAHs στον εσωτερικό αέρα, ιδιαίτερα οι χαμηλού μοριακού βάρους, προέρχονται κυρίως από εσωτερικές πηγές εκπομπής. Υψηλότερες συγκεντρώσεις PAHs παρατηρήθηκαν σε σπίτια με καπνιστές στην Κίνα, την Ιαπωνία και το Ηνωμένο Βασίλειο 
 Σε σπίτια χωρίς καπνιστές, οι απωθητικές ουσίες για έντομα και οι πρακτικές μαγειρέματος ήταν οι κύριες πηγές των 2- και 3-δακτυλίων PAHs αντίστοιχα. Οι PAHs χαμηλού μοριακού βάρους συνδέονταν επίσης με τη χρήση καυστήρων κηροζίνης και τη ρύπανση εξωτερικού αέρα. Οι PAHs υψηλού μοριακού βάρους (5- και 6-δακτυλικοί) στον εσωτερικό αέρα συνδέονταν κυρίως με εξωτερικές πηγές και τα επίπεδά τους ήταν ίσα ή χαμηλότερα από αυτά του εξωτερικού αέρα.


Υπάρχουν αποδείξεις για τη γονοτοξικότητα και την καρκινογένεση πολλών PAHs σε ζωικά είδη, και επιδημιολογικές μελέτες έδειξαν συσχέτιση μεταξύ της έκθεσης σε PAHs και της εμφάνισης καρκίνου σε διάφορους ανθρώπινους ιστούς \cite{ijerph17238972}.


 \subsubsection {Βαρέα Μέταλλα (Toxic Metals)}


Τα βαρέα μέταλλα απελευθερώνονται στην ατμόσφαιρα είτε μέσω ανθρώπινων δραστηριοτήτων είτε μέσω φυσικών διεργασιών . Η εσωτερική ατμοσφαιρική ρύπανση από βαρέα μέταλλα έχει πολλαπλές αιτίες, όπως διείσδυση εξωτερικών ρύπων (σκόνη και έδαφος), κάπνισμα, προϊόντα κατανάλωσης καυσίμων και οικοδομικά υλικά.

Τα βαρέα μέταλλα που περιέχονται στη σκόνη εσωτερικών χώρων μπορούν να εισέλθουν στον ανθρώπινο οργανισμό μέσω εισπνοής, κατάποσης ή επαφής με το δέρμα και να έχουν επιβλαβείς επιπτώσεις στην υγεία. Σύμφωνα με τον Διεθνή Οργανισμό Έρευνας για τον Καρκίνο (IARC),τα βαρέα μέταλλα στον εσωτερικό αέρα ταξινομούνται σε δύο κύριες ομάδες ανάλογα με τις επιδράσεις τους στον άνθρωπο:

\begin{enumerate}
    \item Μη καρκινογόνα στοιχεία, συμπεριλαμβανομένων το κοβάλτιο (Co), το αλουμίνιο (Al), το χαλκό (Cu), το νικέλιο (Ni), το σίδηρο (Fe) και το ψευδάργυρο (Zn)·
    \item Στοιχεία με τόσο καρκινογόνες όσο και μη καρκινογόνες ιδιότητες, στα οποία περιλαμβάνονται το αρσενικό (As), το χρώμιο (Cr), το κάδμιο (Cd) και ο μόλυβδος (Pb). 
\end{enumerate}

Τα πιο κοινά βαρέα μέταλλα,το αρσενικό (As), το χρώμιο (Cr),το κάδμιο (Cd) και ο μόλυβδος Pb, μπορούν να προκαλέσουν καρκίνο , ενώ το κάδμιο και ο μόλυβδος, μαζί με κάποια άλλα, μπορεί να προκαλέσουν και άλλες βλαβερές επιδράσεις, όπως καρδιαγγειακές παθήσεις, καθυστέρηση ανάπτυξης και βλάβες στο νευρικό σύστημα.

Έχει αναφερθεί ότι τα επίπεδα μολύβδου στον εσωτερικό αέρα κυμαίνονται από 5.80 έως 639.10 μg/g, ενώ τα μέγιστα επίπεδα των υπόλοιπων χημικών στοιχείο που αναφέρθηκαν είναι για το αρσενικό 486.80 μg/g, για το αλουμίνιο 7150.00 μg/g, για το χρώμιο 254.00,για το κάδμιο 8.48,για το κοβαλτίο  43.40 μg/g, για το χαλκό 513,00,για το νικέλιο 471,00 μg/g,για τον σίδηρο 4801,00 μg/g και για τον ψευδάργυρο Fe 2293,5 μg/g \cite{ijerph17082927}.
